\documentclass[11pt,letterpaper]{article}
\usepackage[T1]{fontenc}
\usepackage[utf8]{inputenc}
\usepackage{lmodern}
\usepackage[margin=0.9in,headheight=15pt,headsep=18pt]{geometry}
\usepackage{amsmath,amssymb,mathtools,mathrsfs}
\usepackage{microtype}
\usepackage{booktabs,longtable,array}
\usepackage{graphicx}
\usepackage[section]{placeins}
\usepackage{xcolor}
\usepackage{enumitem}
\usepackage{fancyhdr}
\usepackage{titlesec}
\usepackage{tocloft}
\setlength{\cftbeforesecskip}{3pt}
\renewcommand{\cftsecfont}{\small}
\renewcommand{\cftsecpagefont}{\small}
\usepackage{needspace}
\usepackage{xurl}
\usepackage{hyperref}
\hypersetup{hidelinks,pdftitle={Komorebi Large Signal Model: Foundations 0.7.0},pdfauthor={Komorebi P.}}
\setlength{\parindent}{0pt}
\setlength{\parskip}{0.52em}
\setlength{\emergencystretch}{2em}
\setcounter{secnumdepth}{-1}
\setcounter{tocdepth}{1}
\titleformat{\section}{\Large\bfseries}{\thesection}{0em}{}
\titlespacing*{\section}{0pt}{1.7em}{0.7em}
\titleformat{\subsection}{\large\bfseries}{\thesubsection}{0em}{}
\titlespacing*{\subsection}{0pt}{1.15em}{0.5em}
\clubpenalty=10000
\widowpenalty=10000
\displaywidowpenalty=10000
\allowdisplaybreaks[1]
\providecommand{\tightlist}{\setlength{\itemsep}{0pt}\setlength{\parskip}{0pt}}
\providecommand{\passthrough}[1]{#1}
\pagestyle{fancy}
\fancyhf{}
\fancyhead[L]{\small\textsc{Komorebi Large Signal Model}}
\fancyhead[R]{\small Foundations 0.7.0}
\fancyfoot[L]{\small Komorebi P.}
\fancyfoot[R]{\small\thepage}
\renewcommand{\headrulewidth}{0.3pt}
\begin{document}
\hypersetup{pageanchor=false}
\begin{titlepage}
\vspace*{0.5in}
{\small\textsc{Operation Komorebi}\par}
\vspace{0.72in}
{\fontsize{33}{39}\selectfont\bfseries Large Signal Model\par}
\vspace{0.3in}
{\fontsize{21}{27}\selectfont Robust curved certificates\\and noisy geometry\par}
\vspace{0.45in}
\rule{\textwidth}{0.45pt}\par
\vspace{0.3in}
\begin{minipage}{0.95\textwidth}
\large
Propose numerically. Verify exactly.\\Nine statements with proofs.\\Uncertainty and the cost of repair.
\end{minipage}
\vfill
{\Large Komorebi P.\par}
\vspace{0.2in}
{\normalsize Foundations edition 0.7.0\par}
{\normalsize Companion to PR-LASM software release 2.3.0\par}
{\normalsize 9 October 2026\par}
\end{titlepage}
\hypersetup{pageanchor=true}
\begingroup
\footnotesize
\setlength{\cftbeforesecskip}{1pt}
\tableofcontents
\endgroup
\newpage

\section{Part VII. Robust curved certificates and noisy
geometry}\label{part-vii.-robust-curved-certificates-and-noisy-geometry}

The preceding block discovered a curved probability domain from exact
operators. Exact preservation is not the same as stability under
uncertainty. In the worked continuous model, every pure boundary state
can remain on the boundary. A perturbation too small to be conspicuous
in a coefficient table can then move a predicted state outside the
probability domain.

This block separates three operations. A numerical procedure proposes a
geometry. An exact checker decides whether the submitted polynomial
family has a declared margin on that geometry. A separately stated
uncertainty set determines which other families share the guarantee. If
a repair changes the operators, its prediction cost is retained rather
than called an exact coordinate change.

The resulting construction is broader than the preceding exact conformal
search. It accepts nonconformal polynomial families and certifies a full
coefficient box around a repaired estimate. The reference example uses a
deterministic perturbation, not an empirical confidence region.
Statistical coverage is a conditional theorem, not a conclusion of that
example.

Lorentz-positive maps, robust convex constraints and Bernstein
positivity certificates have classical antecedents {[}R8, R13, R21,
R22{]}. The following arguments are self-contained specializations
connecting them to the signal model's representation and evidence
contracts. The nine propositions continue the numbering at 71. Their
count is not a novelty claim or a count of Lean proofs.

\section{80. A margin that has probability
meaning}\label{a-margin-that-has-probability-meaning}

Let \(\mathcal L_d=\{(m,v):m\geq\|v\|_2\}\), with mass row
\(\lambda=(1,0,\ldots,0)\). A mass preserving invertible matrix \(T\)
transports a proposed predictive domain to this cone. Its normalized
base in the original coordinates is

\[
K_T=\{T^{-1}(1,v):\|v\|_2\leq1\}.
\]

Write \(F_a(u)=TB_a(u)T^{-1}\) for the transported operator density.
Observations lie in \([-1,1]\) with reference measure \(du/2\). All
operator densities in this block are polynomial and all equalities
involving integrals use that reference measure. For each action define

\[
\sigma_a(u)=\min_{\|v\|_2\leq1}
\left[(F_a(u)(1,v))_0-\|(F_a(u)(1,v))_{\rm sp}\|_2\right].
\]

The subscript \(\rm sp\) denotes spatial coordinates, not a separate
measured signal. This margin is expressed before normalization. It
simultaneously controls mass orientation and distance from the curved
boundary. It is coordinate dependent, so its uncertainty budget must use
the same \(T\).

\Needspace{5\baselineskip}
\textbf{Proposition 71. Robust oriented margin.} Suppose
\(\sigma_a(u)\geq\kappa_a(u)\) and an uncertain family has transported
errors \(\Delta_a(u)\) satisfying, for every normalized input,

\[
\begin{aligned}
|[\Delta_a(u)(1,v)]_0|&\leq e_{a,0}(u),\\
\|[\Delta_a(u)(1,v)]_{\rm sp}\|_2&\leq e_{a,1}(u).
\end{aligned}
\]

If \(\kappa_a(u)\geq e_{a,0}(u)+e_{a,1}(u)\), every uncertain branch
preserves \(\mathcal L_d\). If the uncertain family has exact integrated
mass and its initial state belongs to \(K_T\), it defines a valid
probability model at every finite horizon. With a strict inequality, the
residual slack is at least \(\kappa_a-e_{a,0}-e_{a,1}\).

\textbf{Proof.} For \(w=F_a(u)(1,v)\) and \(d=\Delta_a(u)(1,v)\),

\[
(w_0+d_0)-\|w_{\rm sp}+d_{\rm sp}\|_2
\geq w_0-\|w_{\rm sp}\|_2-|d_0|-\|d_{\rm sp}\|_2.
\]

The right side is nonnegative under the hypothesis. Homogeneity extends
the assertion from normalized inputs to the entire cone. The only
zero-mass cone element is zero. Integrated mass conservation and
successive cone preservation give nonnegative normalized observation
laws by the same induction as in Proposition 52. A positive-density
branch may be normalized. A zero-density branch is not assigned a
posterior. \(\square\)

This is a uniform deterministic guarantee. The inequalities must cover
the whole observation interval and the whole normalized domain. A noise
scale estimated from a few residuals does not automatically provide
those inequalities. For continuous observations the displayed margin is
a density margin relative to the declared measure, not a probability
assigned to a single observation value.

\section{81. Why exact boundary preservation can have zero noise
radius}\label{why-exact-boundary-preservation-can-have-zero-noise-radius}

An exact certificate can be correct while having no tolerance to
unrestricted perturbations. This is a geometrical fact, not a failure of
the certificate checker.

\Needspace{5\baselineskip}
\textbf{Proposition 72. Boundary fragility under mass preserving
perturbations.} Let \(F\) preserve \(\mathcal L_d\). Suppose a
normalized input \(z=(1,v)\) has a boundary image \(Fz=(m,w)\) with
\(m=\|w\|_2\). For every \(\epsilon>0\), there is a linear perturbation
\(\Delta\) with \(\lambda\Delta=0\) and

\[
\sup_{\|v'\|_2\leq1}\|\Delta(1,v')\|_{\mathcal L_d}=\epsilon
\]

such that \((F+\Delta)z\notin\mathcal L_d\). Here the base norm is
\(\|(m,w)\|_{\mathcal L_d}=\max(|m|,\|w\|_2)\).

\textbf{Proof.} If \(w\ne0\), put \(e=w/\|w\|_2\). If \(w=0\), choose
any spatial unit vector \(e\). Define \(\Delta=\epsilon(0,e)\lambda\).
Its mass row vanishes and its value on every normalized input is
\((0,\epsilon e)\), of base norm \(\epsilon\). If \(w\ne0\), the new
spatial norm is \(m+\epsilon\). If \(w=0\), the original output mass is
zero and the new spatial norm is \(\epsilon\). Either image lies outside
the cone. \(\square\)

For a continuous operator family, a constant spatial perturbation of
this form preserves every integrated mass row and produces the same
obstruction at an observation with a boundary image. It need not
preserve the algebraic restrictions of a narrower uncertainty class.
Perturbing a positive congruence factor and rebuilding its congruence,
for example, is different from perturbing an arbitrary predictive
matrix.

In the preceding exact continuous example, the direct operator at
\(u=0\) is \(I/4\). Every normalized boundary state has a boundary
image. There is therefore no positive unrestricted matrix-error radius
on that domain, even though the original family is exactly valid. A
robust theorem that silently treats its zero slack as positive would be
false.

\section{82. An explicit repair and the price of its
slack}\label{an-explicit-repair-and-the-price-of-its-slack}

In Lorentz coordinates let

\[
N_a(u)=g_a(u)e_0\lambda,\qquad
 g_a(u)\geq g_{\min}>0,\qquad \int g_a(u)\,du/2=1.
\]

This reference instrument reports a state-independent density and resets
to the center. It is a mathematical model modification, not a new
observation. Given an estimate \(\widehat F\), define

\[
\widetilde F_a(u)=(1-\eta)\widehat F_a(u)+\eta N_a(u),
\qquad 0\leq\eta\leq1.
\]

\Needspace{5\baselineskip}
\textbf{Proposition 73. Interior repair with an explicit error budget.}
Suppose a reference family \(F^*\) is valid on \(\mathcal L_d\), and the
combined mass and spatial error of \(\widehat F-F^*\) is at most
\(d_a(u)\) on normalized inputs. Then the margin of
\(\widetilde F_a(u)\) is at least

\[
\boxed{\eta g_a(u)-(1-\eta)d_a(u).}
\]

An additional uncertainty family with combined output error at most
\(\nu_a(u)\) is certified whenever

\[
\eta g_a(u)\geq(1-\eta)d_a(u)+\nu_a(u).
\]

Assume exact integrated mass for the estimate. If its integrated
base-norm discrepancy from \(F^*\) is at most \(\varepsilon\), the
corrected discrepancy is at most

\[
\boxed{(1-\eta)\varepsilon+2\eta.}
\]

When the corrected model is valid on the same cone, with the same
initial state, its length-\(T\) law under any common policy differs from
the reference law by at most

\[
\min\left(1,\frac T2[(1-\eta)\varepsilon+2\eta]\right).
\]

\textbf{Proof.} The margin of \(\widehat F\) is at least \(-d_a\) by
Proposition 71's inequality applied to \(F^*\). Adding the reference
changes only output mass, by \(\eta g_a\), while scaling the remaining
output by \(1-\eta\). This proves the first bound and the robust
implication. Integrated mass is preserved by mixing.

For the error, write

\[
\widetilde F-F^*=(1-\eta)(\widehat F-F^*)+\eta(N-F^*).
\]

On a normalized input, both \(N\) and \(F^*\) distribute total cone mass
one. Their integrated difference has base norm at most two. Integrate
the triangle inequality to obtain the stated budget. The final
implication is the continuous observation-law bound of Proposition 59.
It requires validity of both comparison models on the same cone.
\(\square\)

For constant budgets, a sufficient repair weight is
\((d+\nu)/(g_{\min}+d)\) when this number is at most one. This is not an
optimality claim. The reference example instead fixes a repair weight
and verifies the result independently. If an exact model already has a
robust structured certificate, adding a reset may be unnecessary and
would introduce avoidable bias.

The term \(2\eta\) is deliberately retained. A larger repair creates
more room for uncertainty but can change future predictions more. The
worst-case law bound can become uninformative at long horizons. Neither
positivity nor mixing alone proves long-horizon predictive accuracy.

\section{83. Noisy invariant forms are proposals, not
certificates}\label{noisy-invariant-forms-are-proposals-not-certificates}

The exact identity \(A^{\mathsf T}HA=H\) need not survive estimation
error. A useful proposal can still be obtained from a least-squares
invariant form, but the residual must not be interpreted as a
positive-cone guarantee.

Let \(\mathcal M_A\) map symmetric matrices into a Euclidean product
space by

\[
\mathcal M_A(H)=(A^{\mathsf T}HA-H,\operatorname{tr}H).
\]

Use the Frobenius norm on symmetric matrices and the induced product
norm. Assume \(\mathcal M_{A^*}\) is injective and
\(\mathcal M_{A^*}(H^*)=(0,1)=b\). This fixes the scale of a
one-dimensional invariant-form space.

\Needspace{5\baselineskip}
\textbf{Proposition 74. A posteriori proposal error.} Given
\(\widehat H\), let \(r=\mathcal M_{\widehat A}(\widehat H)-b\). Suppose

\[
\|\mathcal M_{\widehat A}-\mathcal M_{A^*}\|\leq e_M,
\qquad \sigma_{\min}(\mathcal M_{\widehat A})>e_M.
\]

Then

\[
\boxed{
\|\widehat H-H^*\|_F
\leq
\frac{\|r\|+e_M\|\widehat H\|_F}
{\sigma_{\min}(\mathcal M_{\widehat A})-e_M}.
}
\]

One may take \(e_M\) from the bound

\[
\|\mathcal M_{\widehat A}-\mathcal M_{A^*}\|\leq(\|\widehat A\|_2+\|A^*\|_2)\|\widehat A-A^*\|_2.
\]

For fixed centers \(c^*=(I-A^*)^{-1}b^*\) and
\(\widehat c=(I-\widehat A)^{-1}\widehat b\), with
\(\|\widehat A-A^*\|_2\leq e_A\) and \(\|\widehat b-b^*\|_2\leq e_b\),
one has

\[
\|\widehat c-c^*\|_2
\leq\frac{e_b+e_A\|\widehat c\|_2}
{\sigma_{\min}(I-\widehat A)-e_A}
\]

when the denominator is positive.

\textbf{Proof.} The invariant-form difference satisfies

\[
\mathcal M_{A^*}(\widehat H-H^*)
=r-(\mathcal M_{\widehat A}-\mathcal M_{A^*})\widehat H.
\]

The least singular value of \(\mathcal M_{A^*}\) is at least that of
\(\mathcal M_{\widehat A}\) minus \(e_M\). Apply this lower bound and
the triangle inequality. To bound the map perturbation, insert and
subtract \(\widehat A^{\mathsf T}HA^*\) and apply the Frobenius product
inequality. The trace component is unchanged. Finally,

\[
(I-A^*)(\widehat c-c^*)=(\widehat b-b^*)+(\widehat A-A^*)\widehat c.
\]

The same singular-value argument proves the center bound. \(\square\)

The theorem assumes an actual perturbation budget and a corresponding
exact reference problem. The example reports numerical residuals but
does not claim to have statistically certified these budgets. It uses a
trace-normalized least-squares shape, a nearly rank-one image for scale,
and rationalized center and factor. It accepts none of those
computations as a proof of validity. The next checker uses only the
submitted rational geometry and the complete polynomial operators.

A rationalized invertible factor \(L\) specifies \(H=L^{\mathsf T}L\)
exactly. Positive definiteness is then an algebraic fact even when the
factor was suggested by floating-point calculations. It still does not
establish operator invariance.

\section{84. Finite exact checking on a
continuum}\label{finite-exact-checking-on-a-continuum}

For three predictive coordinates, the spatial unit circle admits two
rational charts. For \(s\in\{1,-1\}\) and \(t\in[-1,1]\), put

\[
d(t)=1+t^2,\qquad
r_s(t)=\frac{s}{d(t)}(1-t^2,2t).
\]

The charts cover the two semicircles, including their endpoints. For an
operator \(F_a(u)\) and required margin \(\kappa\geq0\), form polynomial
vectors and scalars

\[
\begin{aligned}
w_{a,s}(u,t)&=F_a(u)(d(t),s(1-t^2),2st)^{\mathsf T},\\
f_{a,s}(u,t)&=w_{a,s,0}(u,t)-\kappa d(t),\\
h_{a,s}(u,t)&=f_{a,s}(u,t)^2-w_{a,s,1}(u,t)^2-w_{a,s,2}(u,t)^2.
\end{aligned}
\]

No square root occurs in these polynomials.

\Needspace{5\baselineskip}
\textbf{Proposition 75. Rational semicircle and Bernstein certificate.}
If \(f_{a,s}\) and \(h_{a,s}\) are nonnegative on \([-1,1]^2\) for both
charts and every action, then \(\sigma_a(u)\geq\kappa\) for every
observation. For rational operators and geometry, nonnegative
tensor-Bernstein coefficients on a complete rational subdivision of that
square give a finite exact witness. If the actual margin is uniformly
strictly greater than \(\kappa\), sufficiently fine dyadic subdivision
succeeds.

\textbf{Proof.} The first polynomial inequality establishes output
orientation after subtracting \(\kappa d\). The second then gives

\[
w_{a,s,0}-\|(w_{a,s,1},w_{a,s,2})\|_2\geq\kappa d.
\]

Divide by \(d>0\) to obtain the required margin on the circle. For fixed
\(u\), the function consisting of affine output mass minus the norm of
affine spatial output is concave in the input. Every point in the closed
disk is a convex combination of circle points. Concavity therefore
extends the lower bound to the disk.

On any rectangle mapped affinely to \([0,1]^2\), the tensor-Bernstein
basis functions are nonnegative and sum to one. A polynomial with
nonnegative coefficients in that basis is nonnegative throughout the
rectangle. If the rectangles cover the square, their inequalities prove
the global statement. The reference verifier checks coverage through a
complete subdivision tree, not by accepting an arbitrary list of
rectangles.

For termination under strict margin, both \(f\) and \(h\) are strictly
positive on the compact square. Their minima are positive. For a
fixed-degree polynomial on a small cell of side at most \(\delta\), its
Taylor expansion at one corner has constant term equal to its value
there and nonconstant coefficients bounded uniformly by \(O(\delta)\).
The linear change from the fixed-degree power basis to the Bernstein
basis is fixed and bounded. Thus every Bernstein coefficient differs
from that constant value by \(O(\delta)\), uniformly over cells. At
sufficiently small \(\delta\) all coefficients are positive. A finite
dyadic subdivision attains that size. \(\square\)

This proof is a specialized use of classical Bernstein range bounds
{[}R22{]}. It does not claim that a negative Bernstein coefficient
proves the polynomial negative. Such a coefficient can merely mean that
further subdivision is needed. A negative evaluated corner, by contrast,
is an actual counterexample. At a zero-margin boundary, finite
termination of this particular witness language is not promised. The
exact conformal and matrix certificates of preceding blocks remain
useful there.

The two tests are both necessary to this sufficient certificate. The
squared inequality alone also admits reversed mass orientation. Both
charts are required. A verifier that omits one chart has not covered the
circle.

\section{85. A finite uncertainty box and the role of
convexity}\label{a-finite-uncertainty-box-and-the-role-of-convexity}

A robust witness should identify its uncertainty family rather than
supply an unexplained tolerance. Suppose the transformed error has
polynomial coefficients

\[
\Delta_a(u)=\sum_{k=0}^{D}u^kE_{a,k},\qquad
|(E_{a,k})_{ij}|\leq r_k.
\]

All bounds are in the submitted Lorentz coordinates. For three
coordinates, normalized inputs have each coordinate of magnitude at most
one.

\Needspace{5\baselineskip}
\textbf{Proposition 76. Coefficient boxes and convex uncertainty.} The
coefficient box satisfies the output bounds of Proposition 71 with

\[
\boxed{e_0=3\sum_{k=0}^{D}r_k,\qquad e_1=6\sum_{k=0}^{D}r_k.}
\]

Only coefficient arrays obeying \(\lambda\int\Delta_a(u)\,du/2=0\) are
included in the normalized uncertainty family. A checked margin at least
\(e_0+e_1\) certifies every such array.

Alternatively, if finitely many complete polynomial operator families
satisfy a common margin \(\kappa\) on the same domain, every convex
combination of those families satisfies that margin. Initial states may
also range over a certified convex hull.

For a fixed geometry, the constraints \(\sigma_a(u)\geq\kappa\), exact
mass, and membership in a convex data-consistency set are convex
constraints on the operator coefficients and initial state. Joint
optimization of geometry and operators is not asserted to be convex.

\textbf{Proof.} Each output scalar is bounded by \(3r_k\) on a
normalized input. Bound the spatial Euclidean norm by the sum of the
absolute values of its two coordinates, obtaining \(6r_k\). Sum over
coefficients using \(|u|\leq1\). Mass conservation is a separate linear
constraint and follows for every family satisfying it.

For convex combinations, output mass is affine and the spatial norm is
convex. Hence the margin of the combined output is at least the same
combination of margins. A convex cone contains convex combinations of
admissible initial states. For fixed geometry and input, the inequality
affine mass minus a norm at least \(\kappa\) is convex in the
coefficients. Intersect over all inputs and observations, then add the
linear mass constraints and the convex data set. \(\square\)

The squared polynomial witness in Proposition 75 is a convenient
verification language. Its coefficient constraints are not automatically
a convex optimization formulation in unknown operators. Convexity of the
underlying admissibility set must not be confused with convexity of
every certificate parameterization.

Likewise, a free uncertainty box in all entries contains families that
violate integrated mass. The affine restriction is essential. Removing
physically valid candidates from a statistical confidence set for some
other reason requires its own coverage argument.

\section{86. What a search can
guarantee}\label{what-a-search-can-guarantee}

The numerical proposal is one practical route. There is also a precise
conditional search theorem showing why a broad class of strict curved
certificates can be found without an exact conformal identity.

\Needspace{5\baselineskip}
\textbf{Proposition 77. Conditional search for strict rational curved
certificates.} Consider finitely many rational polynomial operator
families in three coordinates, each with exact mass, and finitely many
rational normalized initial states. Suppose there exists a real mass
preserving invertible \(T_*\) such that every transformed initial state
lies strictly inside the Lorentz cone and every transformed operator
family has a common positive margin on the observation interval. Then a
rational mass preserving invertible \(T\), a positive rational margin,
and finite rational Bernstein witnesses exist. A dovetailed enumeration
of rational transforms, positive rational margins and subdivision depths
eventually finds such a witness. Whenever it returns, the witness is
sound for the convex hull of the supplied families.

\textbf{Proof.} Inversion and matrix multiplication are continuous near
\(T_*\). Initial membership is strict, and the operator margin is
continuous jointly in the transform, the observation, and the normalized
input. The latter two range over compact sets and there are finitely
many families. Thus the strict inequalities persist in a neighborhood of
\(T_*\). Rational mass preserving matrices are dense in that affine
matrix space, and invertibility persists in a neighborhood. Choose a
rational transform there and a rational positive margin smaller than the
remaining minimum. Proposition 75 supplies finite subdivisions because
there is still a strict gap above the chosen margin. Enumerating all
finite rational descriptions and depths eventually reaches this finite
witness. Proposition 76 proves the convex-hull conclusion. \(\square\)

This is a conditional semidecision result, not a polynomial-time
complexity theorem. It does not cover every valid model. Initial states
on an essential boundary and families with zero margin may violate its
assumptions. The implemented procedure uses a finite rationalized
proposal and bounded subdivision, not this unbounded enumeration.
Failure of that bounded procedure is a withheld certificate rather than
a proof that no other domain works.

The theorem's uncertainty families are supplied in fixed predictive
coordinates. It does not establish that those families cover an unknown
generating process. It demonstrates a route from a strict robust model
sector to finite machine-checkable witnesses.

\section{87. Selection from data does not remove the uncertainty
quantifier}\label{selection-from-data-does-not-remove-the-uncertainty-quantifier}

A domain can be selected after examining fitted operators without a
separate statistical penalty for each geometry tried. That statement
holds only when the final verification ranges over an already justified
confidence set.

\Needspace{5\baselineskip}
\textbf{Proposition 78. Data-selected geometry with honest coverage.}
Suppose a random set \(\mathcal U\) of operators and initial states, in
a declared fixed predictive coordinate system, contains the true model
with probability at least \(1-\delta\). A data-dependent procedure may
propose any transform \(T\) and any domain. If its exact verifier proves
domain membership and invariance for every member of the same set
\(\mathcal U\), then on that coverage event the true model has the
verified domain. No union bound over geometry proposals is needed for
this deterministic implication.

If original-coordinate perturbations satisfy \(\|\Delta\|_2\leq e\),
their combined mass and spatial output error after a known transform
\(T\) is at most

\[
\boxed{2\|T\|_2\|T^{-1}\|_2 e.}
\]

The corresponding transformed uncertainty budget must be used even when
\(T\) was chosen from the data.

\textbf{Proof.} On the coverage event the true model is a member of the
set quantified by the verifier. Its domain properties follow by direct
inclusion, regardless of how \(T\) was selected. For the error, a
normalized Lorentz input has Euclidean norm at most \(\sqrt2\). Thus the
Euclidean norm of \(T\Delta T^{-1}z\) is at most
\(\sqrt2\|T\|_2\|T^{-1}\|_2e\). The sum of its mass magnitude and
spatial norm is at most \(\sqrt2\) times its Euclidean norm. Multiply
the two bounds. \(\square\)

The set must not be shrunk around a convenient candidate without
preserving coverage. Moving the center by a reset mixture changes the
family being certified. A tube around the repaired family need not
contain the original true family. Repeatedly estimating new confidence
sets still requires the stage accounting of Proposition 48. Statistical
rank coverage and observation identifiability remain assumptions of the
upstream procedure.

The final reference certificate explicitly says that no statistical
coverage is claimed. Its uncertainty box is a deterministic, exactly
described set around the repaired operator family. This supports a
reproducible robustness statement without pretending that the example
has estimated an unknown noise distribution.

\section{88. A complete perturbed-operator
construction}\label{a-complete-perturbed-operator-construction}

The input to the new proposer contains only two perturbed polynomial
operators and the initial vector. A separate fixture generator creates
those inputs from the preceding continuous model by adding
\(\pm\delta E(u)\), with opposite signs for the two actions and

\[
\delta=\frac1{2000},\qquad
E(u)=E_0+uE_1+u^2E_2,
\]

where

\[
\begin{aligned}
E_0&=\begin{pmatrix}0&0&0\\-1&2&0\\0&1&-1\end{pmatrix},&
E_1&=\begin{pmatrix}0&1&-1\\1&0&1\\1&-1&0\end{pmatrix},&
E_2&=\begin{pmatrix}0&0&0\\0&-1&1\\1&0&1\end{pmatrix}.
\end{aligned}
\]

The perturbation preserves integrated mass. It is not a statistical
sample. The proposer does not receive the predecessor metric or a hidden
generator.

A noisy operator ratio supplies a center and a least-squares shape. A
finite observation-anchor list supplies a scale proposal. Rationalizing
the center and a triangular metric factor with denominator limit 48
gives

\[
c=\begin{pmatrix}1/2\\1/2\end{pmatrix},\qquad
L=\begin{pmatrix}123/37&-2\\0&109/41\end{pmatrix},
\]

and hence

\[
\boxed{
H=L^{\mathsf T}L=
\begin{pmatrix}15129/1369&-246/37\\-246/37&18605/1681\end{pmatrix}.
}
\]

This is not the exact predecessor metric. The transformed family is not
assumed conformal. The raw fitted family actually fails the proposed
domain: for the direct action at \(u=-1\) and Lorentz input \((1,1,0)\),
its output mass is \(106321/82000>0\), while the output quadratic form
is

\[
-\frac{12009932651317769}{557059220496000000}<0.
\]

This is a genuine rejected input state, not a negative Bernstein
coefficient of uncertain meaning. The normalized output lies outside the
ellipse. No positivity certificate is claimed for the raw fit.

The construction next mixes each operator with the center-reset
reference using \(\eta=1/8\). The exact checker proves a margin of at
least \(1/16\) for the resulting family. It verifies eight polynomial
inequalities across both actions and both semicircle charts. The witness
contains 116 subdivision leaves and 1,780 exactly recomputed nonnegative
Bernstein coefficients. These are arithmetic certificate sizes, not
separate theorems.

The exported family admits mass and spatial perturbation budgets each
equal to \(1/128\), leaving certified slack \(3/64\). In particular, it
covers every degree-two transformed coefficient perturbation with each
entry bounded by \(1/4096\), subject to the exact integrated mass
constraints. Proposition 76 verifies that the whole coefficient box lies
inside the declared output budgets. No uncertainty vertices or future
trajectories are sampled to establish that inclusion.

The certificate is replayed independently from its JSON record. The
record contains rational operators, initial state, transform, margin,
uncertainty definition and complete subdivision trees. Floating-point
proposal diagnostics are stored elsewhere and are not trusted by the
replay routine.

\section{89. The prediction cost needs a common comparison
domain}\label{the-prediction-cost-needs-a-common-comparison-domain}

A successful discovered domain does not imply that the known reference
model has that same domain. To invoke a common-cone law bound, the
example separately checks the corrected family on the predecessor's
exact ellipse. This is a second certificate used only for comparison. It
is not supplied to geometry discovery.

In that comparison domain, the coefficient-wise upper bounds on the raw
fitting discrepancy are

\[
\frac{11}{5000},\qquad \frac{137}{30000},\qquad \frac{37}{8000}.
\]

Using \(\int_{-1}^1|u|^kdu/2=1/(k+1)\) gives the integrated upper budget
\(241/40000\). The same center-reset reference is valid there. Both
models have the same initial state.

\Needspace{5\baselineskip}
\textbf{Proposition 79. Certified repair with a retained comparison
budget.} For the exact construction of Section 88, independent replay
establishes internal validity of the repaired model and its declared
coefficient uncertainty family. Independent comparison-domain replay and
the coefficient bounds above give

\[
\varepsilon_{\rm repaired}
\leq\frac78\frac{241}{40000}+\frac14
=\frac{81687}{320000}.
\]

Consequently, for any common policy and any finite horizon \(T\),

\[
\boxed{
\|P_T-\widetilde P_T\|_{\rm TV}
\leq\min\left(1,\frac{81687T}{640000}\right).
}
\]

This compares the unmodified reference process and the nominal repaired
process. It is not a guarantee that the raw noisy model was valid or
that the repaired uncertainty family contains the original reference.

\textbf{Proof.} The polynomial witnesses and exact mass identities imply
the nominal and robust validity statements by Propositions 71, 75 and
76. The separate comparison witness establishes the required common
domain for the nominal repaired and reference operators. The absolute
coefficient bounds dominate the integrated raw error by the triangle
inequality. Proposition 73 then gives the displayed repair budget and
law bound. All hypotheses are either exact identities in the supplied
witnesses or direct norm bounds from the declared rational coefficients.
\(\square\)

The bound is conservative. It becomes one at horizon eight, so it makes
no nontrivial long-horizon accuracy claim. The example's point is the
separation of a verified all-horizon probability model from a
finite-horizon approximation guarantee. A probability-valid model is not
necessarily a close model.

The noisy proposal, the failed raw fit, the repair and the comparison
bound all appear in the release records. None is replaced by a more
favorable trajectory score. The coefficient box is around the repaired
family in the specified coordinates, not an inferred distribution of
sensor noise.

\section{90. Seventh-block checkpoint}\label{seventh-block-checkpoint}

The cumulative foundation now contains 79 propositions. The new block
extends curved validity checking beyond exact conformal identities to
strict-margin polynomial families, including a concrete
perturbed-operator example. It gives a finite exact certificate over
continuous observations and a continuum of predictive states, a
conditional existence search for strict curved sectors, and the correct
quantifiers for data-selected geometry.

The distinction between noise and model correction is central. A
zero-margin exact model may admit no unrestricted robust neighborhood. A
repaired estimate can have such a neighborhood, but its correction cost
belongs in the prediction budget. A deterministic neighborhood is not a
statistical confidence set until an upstream coverage theorem justifies
that interpretation.

The reference implementation includes a numerical proposal, rationalized
geometry, an exact tensor-Bernstein verifier, a coefficient-box
inclusion check, deliberate rejection controls and an independent
comparison-domain calculation. It does not implement a complete global
optimizer over arbitrary cones, learn empirical coefficient confidence
regions, establish unknown predictive rank, or prove useful long-horizon
accuracy for the worked repair.

The next substantive work is to reduce the conservatism of the repair
and uncertainty descriptions, construct statistically covered operator
sets from realistic dependent signals, and connect these strict
certificates to adaptive nonlinear representation discovery. The current
block establishes a reproducible finite sector rather than declaring
those tasks solved. Production PR-LASM stays at 2.3.0. No new Lean
verification, public publication or license change is claimed.

\begingroup
\small
\raggedright

\section{References and dependencies}\label{references-and-dependencies}

\textbf{{[}R21{]}} Roland Hildebrand. \emph{An LMI description for the
cone of Lorentz-positive maps}. Author manuscript, 17 January 2006.
Primary manuscript inspected, especially the introduction and
definitions on pages 1--2. It supplies classical context for
Lorentz-positive maps and the existence of convex descriptions. The
present margin and semicircle certificate arguments are proved directly
and do not claim to implement the manuscript's full matrix-inequality
construction. Source:
\path{https://optimization-online.org/wp-content/uploads/2005/12/1260.pdf}.

\textbf{{[}R22{]}} Tareq Hamadneh, Hassan Al-Zoubi, Hamza Alzaareer and
Rafael Wisniewski. \emph{Optimization and Positivity Certificates of
Rational Functions using Bernstein Form}. arXiv:1906.11037, 2019.
Primary manuscript inspected, including its Bernstein range enclosure
discussion on printed pages 3--4. It supplies established
positivity-certificate context. Proposition 75 gives its own
tensor-product rectangle argument and strict-subdivision proof. Source:
\path{https://arxiv.org/abs/1906.11037}.

\textbf{Inherited dependencies.} Proposition 52 defines continuous
instrument semantics. Proposition 59 supplies the common-domain
continuous-law bound. Propositions 62--63 distinguish exact ellipsoid
proposals from verification. Proposition 48 supplies repeated
confidence-set accounting. The new block restates the specific repair
and uncertainty inequalities it uses. The unchanged 0.6.0 archive
contains the full predecessor and all earlier references.

\endgroup
\end{document}
